% 调和函数（综述）
% license CCBYSA3
% type Wiki

本文根据 CC-BY-SA 协议转载翻译自维基百科\href{https://en.wikipedia.org/wiki/Harmonic_function}{相关文章}。


在数学、数理物理以及随机过程理论中，调和函数（harmonic function）是指满足拉普拉斯方程（Laplace's equation）的函数。  
设
\[
f : U \to \mathbb{R}~
\]
其中 \(U \subseteq \mathbb{R}^n\) 为一个开集，且 \(f\) 是定义在 \(U\) 上的二阶连续可微函数（即 \(f \in C^2(U)\)）。若 \(f\) 在 \(U\) 上处处满足
\[
\frac{\partial^2 f}{\partial x_1^2}
+ \frac{\partial^2 f}{\partial x_2^2}
+ \cdots
+ \frac{\partial^2 f}{\partial x_n^2}
= 0~
\]
则称 \(f\) 为一个\textbf{调和函数}。  
该方程通常记作
\[
\nabla^2 f = 0,
\quad \text{或等价地} \quad
\Delta f = 0.~
\]
\subsection{“Harmonic” 一词的词源}
“调和函数（harmonic function）”中的形容词“harmonic”（意为“谐和的”或“调和的”）起源于一根被拉紧的弦上某一点所经历的谐振运动（harmonic motion）。描述此类运动的微分方程的解可以用正弦函数（sine）与余弦函数（cosine）来表示，因此这些函数被称为谐波（harmonics）。傅里叶分析（Fourier analysis）研究如何将单位圆上的函数展开为这些谐波的级数。若考虑高维空间中单位 \(n\) 球面上的谐波类比，便得到所谓的球谐函数（spherical harmonics）。这些函数满足拉普拉斯方程。随着时间的推移，“harmonic”一词被推广，用来指代所有满足拉普拉斯方程的函数。\(^\text{[1]}\)  

